% !TEX root = ../report.tex
\chapter{Исследовательская часть}
\section{Методика сравнения}
Временная эффективность оценивается числом выполненных сравнений искомого значения с элементами массива. Для каждого фиксированного размера лучший случай означает минимум, а худший --- максимум по всем возможным запросам, включая отсутствие значения. Сравниваются строго возрастающие массивы. Все формулы ниже относятся к $N\geq1$. При $N=0$ число сравнений равно нулю. Записи $\lfloor z\rfloor$ и $\lceil z\rceil$ означают округление числа $z$ вниз и вверх до целого соответственно. Например, $\lfloor5/2\rfloor=2$, а $\lceil5/2\rceil=3$.

Эксперимент выполнен в среде Fedora Linux 44, на процессоре AMD Ryzen 5 5600H с 12 логическими процессорами. Доступный операционной системе объём оперативной памяти составляет около 14 ГиБ. Стенд подсчёта сравнений собран GCC 16.1.1 без оптимизации, с параметром \texttt{-O0}.

\section{Число сравнений линейного поиска}
При наличии значения первое равенство на позиции $i$ требует $i+1$ сравнений. Если значение отсутствует, проверяются все $N$ элементов. Поэтому для непустого массива
\begin{samepage}
\begin{equation}\label{eq:linear}
 C_{\mathrm{lin}}^{\min}(N)=1,\qquad C_{\mathrm{lin}}^{\max}(N)=N.
\end{equation}
\indent где $C_{\mathrm{lin}}^{\min}$ и $C_{\mathrm{lin}}^{\max}$ --- минимальное и максимальное число сравнений линейного поиска.
\end{samepage}

Лучший случай наступает при $x=a_0$. Худший --- при отсутствии $x$ либо при первом появлении $x$ на последней позиции. Асимптотические оценки составляют $\Theta(1)$ и $\Theta(N)$ соответственно.

\section{Число сравнений классического бинарного поиска}
Каждая итерация содержит сравнение $a_M<x$. Если средний элемент не меньше искомого значения, проверяется $a_M>x$. Если искомое значение равно среднему элементу исходного массива, выполняется два сравнения. Для $N=1$ и $N=2$ поиск значения, большего всех элементов массива, заканчивается после одной итерации и одного сравнения. Для $N\geq3$ результат не может определиться после одного сравнения. Получаем
\begin{samepage}
\begin{equation}\label{eq:classic}
 C_{\mathrm{cl}}^{\min}(N)=\begin{cases}1,&N\leq2,\\2,&N\geq3,\end{cases}
 \qquad C_{\mathrm{cl}}^{\max}(N)=2\lfloor\log_2N\rfloor+2.
\end{equation}
\indent где $C_{\mathrm{cl}}^{\min}$ и $C_{\mathrm{cl}}^{\max}$ --- минимум и максимум числа сравнений классического бинарного поиска.
\end{samepage}

Худший случай достигается при $x<a_0$: выбирается левая половина длины $\lfloor N/2\rfloor$, и на каждой непустой итерации выполняются два сравнения. Так как число элементов на каждой итерации уменьшается примерно вдвое, для массива из $N$ элементов потребуется не более $\lfloor\log_2N\rfloor+1$ итераций. В лучшем случае оценка составляет $\Theta(1)$, в худшем --- $\Theta(\log N)$.

\section{Число сравнений бинарного поиска без ветвления}
В основном цикле шаг уменьшается вдвое, и на каждой итерации выполняется одно сравнение. После цикла сравниваются элемент с индексом \texttt{begin} и искомое значение. Если полученная позиция находится внутри массива, выполняется ещё одна проверка на равенство. Подготовка поиска также может потребовать одного сравнения, если размер массива не является степенью двойки.

Если $N$ является степенью двойки, цикл выполняет $\log_2N$ сравнений. Поэтому лучший случай требует $\log_2N+1$ сравнений, а худший --- $\log_2N+2$. Лучший случай наступает, когда искомое значение больше всех элементов: полученная позиция находится за концом массива и проверка равенства не выполняется.

Для остальных размеров число итераций зависит от того, какая часть массива выбрана перед циклом. Обозначим через $p=\texttt{bit\_floor}(N)$ наибольшую степень двойки, не превышающую $N$. Худший случай достигается при значении, меньшем всех элементов: алгоритм остаётся в начале массива. Лучший --- при значении, большем всех элементов: выбирается конечная часть массива. Если $N=p+1$, поиск завершается сразу после первого сравнения. Иначе в конечной части остаётся $N-p-1$ элементов, и шаг задаётся функцией \texttt{bit\_ceil}.

Таким образом, для непустого массива получаем:
\begin{samepage}
\begin{equation}\label{eq:branchcount}
\begin{aligned}
 C_{\mathrm{bl}}^{\min}(N)&=\begin{cases}
 \log_2N+1,&N=p,\\
 1,&N=p+1,\\
 \lceil\log_2(N-p-1)\rceil+2,&N>p+1,
 \end{cases}\\[2mm]
 C_{\mathrm{bl}}^{\max}(N)&=\begin{cases}
 \log_2N+2,&N=p,\\
 \lfloor\log_2N\rfloor+3,&N>p.
 \end{cases}
\end{aligned}
\end{equation}
\indent где $C_{\mathrm{bl}}^{\min}$ и $C_{\mathrm{bl}}^{\max}$ --- минимальное и максимальное число сравнений бинарного поиска без ветвления;\par
\noindent $p$ --- наибольшая степень двойки, не превышающая $N$.
\end{samepage}

В худшем случае число сравнений имеет оценку $\Theta(\log N)$. Лучший случай зависит от размера: при $N=p+1$ достаточно одного сравнения, а при размерах, равных степени двойки, число сравнений растёт логарифмически. Проверки и вычисления внутри \texttt{bit\_floor} и \texttt{bit\_ceil} не сравнивают искомое значение с элементами массива и в выбранный счётчик не входят.

\section{Число сравнений рекурсивного бинарного поиска}
Рекурсивная реализация начинает проверку непустого подмассива с равенства. Поэтому лучший случай, $x=a_{\lfloor N/2\rfloor}$, требует одного сравнения. Если средний элемент не равен искомому значению, выполняются два сравнения: проверка равенства и проверка \texttt{a[M] < x}. В худшем случае таких вызовов будет $\lfloor\log_2N\rfloor+1$, после чего функция получит пустую часть массива и завершится. Следовательно,
\begin{samepage}
\begin{equation}\label{eq:reccount}
 C_{\mathrm{rec}}^{\min}(N)=1,\qquad C_{\mathrm{rec}}^{\max}(N)=2\lfloor\log_2N\rfloor+2.
\end{equation}
\indent где $C_{\mathrm{rec}}^{\min}$ и $C_{\mathrm{rec}}^{\max}$ --- минимальное и максимальное число сравнений рекурсивного бинарного поиска.
\end{samepage}

Если значение найдено в $k$-м вызове, до этого выполнено $k-1$ вызовов с двумя сравнениями, а последний требует одного сравнения. Всего получается $2k-1$. Если значение не найдено, каждый из $k$ вызовов с непустой частью массива требует двух сравнений, всего --- $2k$. Асимптотически лучший случай имеет оценку $\Theta(1)$, а худший --- $\Theta(\log N)$.

\section{Дополнительная память}
Память входного массива одинакова для всех алгоритмов и составляет $N\cdot\mathrm{sizeof}(\texttt{int})$ байт. Она не включается в дополнительную память поиска. Линейный алгоритм хранит индекс; нерекурсивные бинарные алгоритмы хранят ограниченное число индексов, границ и временных значений. Их дополнительная память имеет оценку $\Theta(1)$.

При рекурсивном поиске в стеке сохраняются данные каждого активного вызова: параметры, локальные переменные и адрес возврата. Поскольку каждый вызов порождает не более одного следующего вызова, дерево рекурсивных вызовов представляет собой цепочку. Её высота равна числу одновременно активных вызовов. Эти данные образуют кадр вызова, размер которого не зависит от длины массива. В худшем случае после непустых участков вызывается функция для пустого участка. Поэтому
\begin{samepage}
\begin{equation}\label{eq:memory}
 H_{\max}(N)=\lfloor\log_2N\rfloor+2.
\end{equation}
\indent где $H_{\max}$ --- максимальное число одновременно активных вызовов, включая вызов для пустого участка.
\end{samepage}

В лучшем случае выполняется один вызов и требуется $\Theta(1)$ дополнительной памяти. В худшем случае число кадров растёт как $\log N$, поэтому расход памяти составляет $\Theta(\log N)$. Точный объём в байтах зависит от компилятора и архитектуры. Оценка относится к выполнению рекурсивных вызовов без их устранения оптимизатором.

Стенд дополнительно хранит отсортированную копию, что требует $\Theta(N)$ памяти. Это общая подготовка данных для бинарных алгоритмов, а не дополнительная память их функций. Пузырьковая сортировка в используемой программе выполняет $N(N-1)/2$ сравнений элементов и имеет трудоёмкость $\Theta(N^2)$; исходный порядок сохраняется. Затраты сортировки не включаются в таблицы стоимости одного поискового запроса. При изменении массива упорядоченность копии необходимо восстанавливать.

\section{Экспериментальные результаты}
Для каждого размера создан массив $a_i=2i$, $0\leq i<N$. Поиск выполняется для каждого целого значения $x$ от $-1$ до $2N$ включительно. Такой набор охватывает присутствующие значения, каждый промежуток между соседними элементами и оба внешних диапазона. Поскольку поведение алгоритмов определяется отношениями порядка и равенства, этот перебор даёт точные минимум и максимум для выбранного массива, а не случайную выборку.

Результаты представлены в таблице~\ref{tab:counts}. В каждой ячейке указаны минимум и максимум числа сравнений, разделённые косой чертой. Для пустого массива все четыре функции выполняют ноль сравнений элементов с $x$.
\newcommand{\comparisonrows}{%
1 & 1 / 1 & 1 / 2 & 1 / 2 & 1 / 2\\\hline
2 & 1 / 2 & 1 / 4 & 2 / 3 & 1 / 4\\\hline
3 & 1 / 3 & 2 / 4 & 1 / 4 & 1 / 4\\\hline
8 & 1 / 8 & 2 / 8 & 4 / 5 & 1 / 8\\\hline
10 & 1 / 10 & 2 / 8 & 2 / 6 & 1 / 8\\\hline
16 & 1 / 16 & 2 / 10 & 5 / 6 & 1 / 10\\\hline
32 & 1 / 32 & 2 / 12 & 6 / 7 & 1 / 12\\\hline
64 & 1 / 64 & 2 / 14 & 7 / 8 & 1 / 14\\\hline
128 & 1 / 128 & 2 / 16 & 8 / 9 & 1 / 16\\\hline
256 & 1 / 256 & 2 / 18 & 9 / 10 & 1 / 18\\\hline
1024 & 1 / 1024 & 2 / 22 & 11 / 12 & 1 / 22\\\hline%
}

\begin{table}[H]
\caption{Экспериментальные минимум и максимум числа сравнений}\label{tab:counts}
\centering\tabfont
\begin{tabularx}{\textwidth}{|r|Y|Y|Y|Y|}
\hline
$N$ & Линейный & Классический бинарный & Бинарный без ветвления & Рекурсивный бинарный\\\hline
\comparisonrows
\end{tabularx}
\end{table}


\section{Сопоставление алгоритмов}
Итоговое сравнение приведено в таблице~\ref{tab:summary}. Формулы справедливы для $N\geq1$; условия относятся к строго возрастающему массиву. Дополнительная память указана для функции поиска без входного массива и отсортированной копии. В таблице $m=\lfloor N/2\rfloor$ --- индекс среднего элемента. Число сравнений алгоритма без ветвления определяется формулой~\eqref{eq:branchcount}.
\begin{table}[H]
\caption{Сравнение ресурсной эффективности алгоритмов поиска}\label{tab:summary}
\centering\tabfont
\begin{tabularx}{\textwidth}{|P{27mm}|P{29mm}|P{27mm}|Y|P{22mm}|}
\hline
Алгоритм & Лучший случай и условие & Худший случай и условие & Ограничения & Память, худший случай\\\hline
Линейный & $1$; $x=a_0$ & $N$; $x$ отсутствует или впервые найден на позиции $N-1$ & Нет требования упорядоченности & $\Theta(1)$\\\hline
Классический бинарный & $1$ при $N\leq2$, $x>a_{N-1}$; $2$ при $N\geq3$, $x=a_m$ & $2\lfloor\log_2N\rfloor+2$; $x<a_0$ & Массив по неубыванию & $\Theta(1)$\\\hline
Бинарный без ветвления & $C_{\mathrm{bl}}^{\min}(N)$; $x>a_{N-1}$ & $C_{\mathrm{bl}}^{\max}(N)$; $x<a_0$ & Массив по неубыванию & $\Theta(1)$\\\hline
Рекурсивный бинарный & $1$; $x=a_m$ & $2\lfloor\log_2N\rfloor+2$; $x<a_0$ & Массив по неубыванию & $\Theta(\log N)$\\\hline
\end{tabularx}
\end{table}

При $N=1024$ худший случай линейного поиска требует 1024 сравнений, классического и рекурсивного бинарных --- 22, алгоритма без ветвления --- 12. Однако для запроса, равного среднему элементу, рекурсивный алгоритм выполняет одно сравнение, классический --- два, а алгоритм без ветвления --- 12. Поэтому предпочтительность алгоритма зависит от сценария запросов.

\section*{Вывод}
Для всех исследованных размеров число сравнений совпало с теоретическими оценками. В худшем случае бинарные алгоритмы требуют логарифмического числа сравнений, линейный --- $N$. На исследованных размерах $N\geq8$ алгоритм без ветвления выполняет меньше сравнений в худшем случае, чем классический и рекурсивный бинарные алгоритмы. Однако при поиске среднего элемента он не завершает цикл сразу после обнаружения равенства. Рекурсивный вариант имеет логарифмические затраты памяти. При выборе бинарного поиска нужно также учитывать предварительную сортировку массива.
